\track{T5\quad 线算子、特征膜、单极子与 Hecke}{本轨道 21 条：物理侧的 Wilson/'t Hooft 线算子与磁单极子，就是数学侧的 Hecke 算子与 Hecke 特征层。这是全书最硬、也最漂亮的一一对应。}

\par\medskip\noindent\textbf{第一组（E）：线算子、特征膜与 Hecke 算子}\par

\begin{word}{Wilson 算子}{Wilson operator}{wilson-op}{\S6.1, p.~79}
\mathdef{设 $M$ 是光滑四维流形，$G$ 是紧李群，$A$ 是 $M$ 上 $G$-联络，$R$ 是 $G$ 的有限维复表示。沿曲线 $S\subset M$ 的平行移动为 $P\exp\int_S A\in G$，其 $R$-迹 \[ W_R(S)=\operatorname{Tr}_R P\exp\int_S A \] 称为 Wilson 线。在扭化理论中，为保持超对称须把某个标量场并入联络，取复联络 $\mathcal{A}=A+\mathrm{i}\varphi$ 的和乐，此时 $W_R(S)$ 才与拓扑超荷相容。}
\mathex{取 $G=SU(2)$、$C=\mathbb{P}^1$、$R$ 为二维基本表示，$S$ 为 $M=\mathbb{R}\times C$ 中沿时间方向的直线 $\mathbb{R}\times\{x\}$，则 $W_2(S)$ 取值在 $2\times 2$ 矩阵中；若把时间方向取成圆 $S^1$ 并取迹，就得到取值在 $\mathbb{C}$ 中的 Wilson 圈。$S^1$ 与 $C$ 上一点 $x$ 的组合给出只依赖 $(R,x)$ 的算子。}
\physdef{物理上 Wilson 算子是有序算子（order operator）：先写下一个经典场的和乐泛函，再把它量子化。它携带一个外源电荷，期望值测量沿 $S$ 的电通量。其超对称版本只在扭化参数 $\Psi=\infty$ 一侧存在；在一般的 $\Psi$ 下它只是 BPS 的，关联函数依赖度量。}
\physex{在纯 $U(1)$ 规范理论中，电荷 $e$ 的 Wilson 线 $W(e,S)$ 与磁荷 $m$ 的 't Hooft 线 $T(m,S')$ 满足 \[ W(e,S)\,T(m,S')=\mathrm{e}^{2\pi\mathrm{i}em\,\operatorname{Link}(S,S')}\,T(m,S')\,W(e,S), \] 其中 $\operatorname{Link}$ 是环绕数。这个相位是线算子代数的基本不变量；在 $SU(2)$ 理论中它还决定基本表示与伴随表示 Wilson 线的 OPE。}
\bridge{对应规则：Wilson 算子 $W_R(S)$ 对应 Hecke 算子 $H_{R,x}$，其中 $x$ 是 $S$ 与黎曼面 $C$ 的交点，$R$ 是 ${}^LG$ 的表示（在电侧即 $G$ 的表示）。自由参数是表示 $R$ 与点 $x\in C$；在膜的层面，Wilson 算子的作用是「与泛丛的表示版本 $\mathcal{E}_p(R)$ 作张量积」。原文见 \S6.1 p.~79、\S8（p.~100 起）与 \S11。}
\intuit{一句话：Wilson 线就是「把电荷沿一条线拖着走」的和乐，而 Hecke 算子正是「沿 $C$ 上一点拖动丛」的操作；两者都只依赖一个表示与一个点，所以在扭化理论里是同一件事的两面。}
\pitfall{按字面理解会以为 Wilson 算子是某个固定 Hilbert 空间上的普通算子；其实它是线算子，开线给出矩阵、闭圈取迹才给出数，其 OPE 系数是向量空间而不是数（只有把时间取成圆并取迹，才退化为维数）。另一个误解是以为其超对称版本对任意 $\Psi$ 都存在。}
\end{word}

\begin{word}{'t Hooft 算子}{'t Hooft operator}{thooft-op}{\S6.2, p.~81}
\mathdef{设 $S\subset M$ 是嵌入曲线，$\rho:U(1)\to G$ 是同态。沿 $S$ 挖去半径 $\varepsilon$ 的管，并要求管壁上的场具有由 $\rho$ 给出的奇异性：对 $G=U(1)$ 取 \[ F\sim\frac{\mathrm{i}}{2}\star_3\mathrm{d}\Big(\frac{1}{r}\Big),\qquad \varphi\sim\frac{\mathrm{i}}{2r}\mathrm{d}s, \] 然后在满足该奇异性的场空间上作路径积分，就得到 't Hooft 算子 $T(\rho,S)$。它只依赖 $\rho$ 的共轭类；共轭类由 ${}^LG$ 的最高权分类，故记作 $T({}^LR,S)$，其中 ${}^LR$ 是 ${}^LG$ 的不可约表示。}
\mathex{$G=U(1)$：磁荷 $m$ 的算子取 $F\sim(\mathrm{i}m/2)\star_3\mathrm{d}(1/r)$，绕 $S$ 的磁通量为 $m\in\mathbb{Z}$，等价于 $\mathbb{P}^1$ 上的 $\mathcal{O}(m)$ 型跃迁 $z^{-m}$。$G=SU(2)$：$\rho(\mathrm{e}^{\mathrm{i}\alpha})=\mathrm{diag}(\mathrm{e}^{\mathrm{i}n\alpha/2},\mathrm{e}^{-\mathrm{i}n\alpha/2})$，$n\ge0$，$n=1$ 为基本磁荷；标签 $n$ 同时是 ${}^LG=SO(3)$ 的 $(n+1)$ 维表示的最高权。}
\physdef{物理上 't Hooft 算子是无序算子：不规定和乐取值，而规定场沿 $S$ 的奇异性，在满足奇异性的场空间上作路径积分；它是「把磁荷钉在一条线上的奇异源」。这个定义需要 $S$ 的法丛定向（因为用到 $\star$），纯 't Hooft 算子只在 $\Psi=0$ 时是拓扑算子。}
\physex{四维 $U(1)$ 理论中 $T(m,S)$ 就是插入一条狄拉克单极子线的算子；$SU(2)$ 理论中 $T(1,S)$ 附近的奇异解就是 \S6.2 讨论的磁单极子。't Hooft 算子的 OPE 系数由 Hecke 修正空间的相交上同调给出，最终等于 ${}^LG$ 的表示环（\S10.4）。}
\bridge{核心对应：${}^LG$ 的最高权 ${}^Lw$（GNO 磁荷）$\leftrightarrow$ 仿射 Grassmann 流形中的点 $\mathrm{diag}(z^{-m_1},\dots,z^{-m_N})$ $\leftrightarrow$ Hecke 修正的型；同一个 $\rho$ 的这几次化身见原文式 (9.32)（\S9.2, p.~147）。自由参数是 ${}^Lw$ 与点 $x\in C$；原文见 \S6.2 p.~81、\S9.1--9.2、\S10。}
\intuit{一句话：'t Hooft 线的奇异性就是「允许在一点把丛改写成什么样」的指令，而 Hecke 修正正是按指令改写丛，所以 't Hooft 算子就是 Hecke 算子。}
\pitfall{按字面理解会把 't Hooft 算子当成磁单极子：前者是支撑在曲线上的线算子，后者是算子的奇异解（或态）。也会误以为它和 Wilson 算子是同一理论里可直接相乘的两个算子；实际上它们只在 S-对偶下互换，且只有有理 $\Psi$ 时才能拼成复合算子 $WT_{n,m}$，纯 't Hooft 算子只在 $\Psi=0$ 存在。}
\end{word}

\begin{word}{线算子}{line operator}{line-op}{\S6, p.~79}
\mathdef{线算子是支撑在一维子流形上的算子。给定流形 $M$ 上的量子场论，线算子由三项数据给出：一条带边界的嵌入曲线 $S\subset M$、沿 $S$ 的经典场泛函或奇异性条件、以及在端点处作用于边界条件的向量空间。数学上，线算子构成一个范畴 $\mathcal{L}$：对象是线算子，态射是端点处的界面；线算子作用在膜（边界条件）上给出函子。}
\mathex{$G=SU(2)$、$M=\mathbb{R}\times\mathbb{P}^1$：沿 $\mathbb{R}\times\{x\}$ 的基本表示 Wilson 线 $W_2$ 与伴随表示 Wilson 线 $W_3$ 的乘积含 $W_2$ 项，多重数空间为 $\operatorname{Hom}_{SU(2)}(\mathbf{2}\otimes\mathbf{3},\mathbf{2})\cong\mathbb{C}$；而 $W_2\cdot W_2$ 不含 $W_2$ 项，只含 $W_3$ 与 $W_1$（因为 $\mathbf{2}\otimes\mathbf{2}=\mathbf{3}\oplus\mathbf{1}$）。}
\physdef{与局域算子相对：局域算子由 0 维支撑上的局部场表达，线算子必须指定沿曲线的外源或奇异性，因而不是局部泛函。在拓扑场论中线算子是基本的非局域可观测量，其作用完全由闭圈上的期望值与 OPE 决定。}
\physex{四维 $\mathcal{N}=4$ 扭化理论的两大类线算子是 Wilson 线（有序）与 't Hooft 线（无序）。把理论紧化到 $\Sigma\times C$ 后，支撑在 $\Sigma$ 上的线算子成为二维理论的边界条件（膜），而支撑在 $\Sigma$ 上一点、沿 $C$ 方向的线算子成为二维理论中的点算子。}
\bridge{对应规则：线算子的作用范畴与 Hecke 算子的作用范畴同构。Wilson 线对应沿点的 Hecke 算子（张量积实现），'t Hooft 线对应沿对偶点的 Hecke 算子（修正实现），一般线算子对应 Hecke 特征层。原文见 \S6 p.~79 与 \S9--\S10。}
\intuit{一句话：线算子就是「延展的荷」，Hecke 算子是「延展的改写」，两者都靠一条曲线上的一点来定位，所以词典把它们对成同一类对象。}
\pitfall{按字面理解会以为开线与闭圈是同一套代数：开线的端点是边界条件、给出矩阵，闭圈取迹才给出数，二者 OPE 结构不同。也会误以为线算子只由其在闭圈上的期望值决定，但不同的线算子可以有相同的期望值而不同的 OPE。}
\end{word}

\begin{word}{圈算子}{loop operator}{loop-op}{\S6.1, p.~79}
\mathdef{圈算子是支撑在闭曲线 $S$ 上的线算子。对闭曲线 $S$，若线算子由 $G$-联络的和乐给出并取表示 $R$ 的迹，则得规范不变函数 $W_R(S)=\operatorname{Tr}_R P\exp\oint_S A$。一般地，圈算子由「和乐加一个与 $R$ 相关的端点粘合」给出；数学上它是线算子范畴中闭态射的迹（范畴迹），因此只依赖线算子的同构类。}
\mathex{$G=U(1)$：圈算子为 $W(e,S)=\mathrm{e}^{\mathrm{i}e\oint_S A}$，是规范不变的 $U(1)$ 相位；若 $S$ 与磁荷 $m$ 的 't Hooft 圈 $T(m,S')$ 环绕，则相位为 $\mathrm{e}^{2\pi\mathrm{i}em}$。$G=SU(2)$：$W_2(S)$ 是 $2\times 2$ 和乐的迹，展开后含 $\operatorname{Tr}(F_{z\bar z})^2$ 等局域项的积分。}
\physdef{圈算子的物理意义是「测量沿闭路的荷」；它的期望值是可观测量的基本数据，Wilson 圈的期望值在 $\mathcal{N}=4$ 理论中可由局域化精确计算。取时间方向为圈时，圈算子等价于对 Hilbert 空间取迹（配分函数），这也是其 OPE 系数为整数的原因。}
\physex{$SU(2)$ 理论中闭 Wilson 圈 $W_2(S)$ 的 OPE 系数：取迹后 $\mathbf{2}\otimes\mathbf{2}=\mathbf{3}\oplus\mathbf{1}$ 的系数退化为维数 $1,1$；而开（静态）Wilson 线的 OPE 系数是向量空间 $\operatorname{Hom}(\mathbf{2}\otimes\mathbf{2},\mathbf{3})\cong\mathbb{C}$。原文 \S10.4 正是用这个对比说明「取迹把向量空间换成了维数」。}
\bridge{对应规则：圈算子对应 Hecke 特征层的「特征值」一侧——在闭圈上取迹相当于在 Hecke 特征层上取 $V_E$ 的纤维维数。数学侧 Hecke 特征层满足 $H_V(\mathcal{F})\simeq V_E\boxtimes\mathcal{F}$，其闭圈版本就是取维数。原文见 \S6.1 p.~79、\S10.4 p.~169。}
\intuit{一句话：圈算子是线算子取迹后的数，就像 Hecke 算子在特征层上作用后只剩下一个向量空间作为「特征值」；两者都是把矩阵压成一个数（或一个层）的操作。}
\pitfall{按字面理解会以为圈算子只是线算子的一种、并不重要；其实正是取迹才使圈算子成为规范不变的观测量，也才使 OPE 系数从向量空间退化为整数。也不要把圈算子与局部算子混为一谈：前者支撑在整条闭路上，后者只支撑在一点。}
\end{word}

\begin{word}{有序算子}{order operator}{order-op}{\S6.2, p.~81}
\mathdef{设理论的经典场空间为 $\mathcal{F}$，路径积分为 $\int_{\mathcal{F}}\mathrm{e}^{\mathrm{i}S}$。有序算子是由经典场的泛函 $\mathcal{O}:\mathcal{F}\to\mathbb{C}$ 给出的算子：插入 $\mathcal{O}$ 就是把该泛函乘进路径积分的被积函数。数学上这是场空间上函数代数的乘法算子（几何量子化中的 Toeplitz 型算子）；其定义不需要预先指定任何奇异性，也不需要挖去曲线的邻域。}
\mathex{$G=U(1)$、$M=\mathbb{R}^4$：$\mathcal{O}=F_{\mu\nu}(x)$ 是有序算子；沿曲线 $S$ 的 Wilson 线 $W(e,S)=\mathrm{e}^{\mathrm{i}e\oint_S A}$ 也是有序算子，因为它就是经典和乐的一个泛函。$SU(2)$ 情形同理，只需把 $\mathrm{e}^{\mathrm{i}e\oint A}$ 换成表示 $R$ 的迹。}
\physdef{有序算子（order operator）与无序算子相对：有序算子用经典场的表达式定义、再量子化，测量的是场的局部或延展值；无序算子则通过规定奇异性、改变被积场空间来定义。Wilson 算子是有序算子的典型例子，'t Hooft 算子是无序算子的典型例子。}
\physex{在 $SU(2)$ 规范理论中，基本表示 Wilson 线 $W_2(S)$ 是有序算子；它携带外源电荷，测量沿 $S$ 的电通量。与其相对，'t Hooft 线 $T(1,S)$ 是无序算子，测量绕 $S$ 的磁通量；两者在 $\mathbb{R}^4$ 中互不交换，其交换相位由环绕数给出。}
\bridge{对应规则：有序算子对应「张量积型」的 Hecke 作用（Wilson 侧），无序算子对应「修正型」的 Hecke 作用（'t Hooft 侧）；同一个 Hecke 算子的两种实现由 S-对偶联系。原文见 \S6.2 p.~81 与 \S9。}
\intuit{一句话：有序算子就是「把经典量直接搬进量子理论」，所以它对应的数学操作是乘法与张量积；无序算子改变的是场的定义域，所以对应的是改写丛的 Hecke 修正。}
\pitfall{按字面理解会以为「有序/无序」说的是算符顺序（normal ordering）；其实它指算子是由经典泛函（order）定义还是由奇异性（disorder）定义，与正规序无关。}
\end{word}

\begin{word}{无序算子}{disorder operator}{disorder-op}{\S6.2, p.~81}
\mathdef{设 $S\subset M$ 是余维 2 的子流形，并给定 $S$ 法丛的一个定向与同态 $\rho:U(1)\to G$。无序算子是如下路径积分：在 $M\setminus S$ 的场空间上积分，并要求绕 $S$ 的小圈上的和乐共轭于 $\rho(\mathrm{e}^{\mathrm{i}\theta})$。等价地，把 $S$ 挖去细管并在管壁上给定边界条件。数学上它改变的是场的定义域与奇异性，而不是乘一个泛函。}
\mathex{$G=U(1)$：$\rho(\mathrm{e}^{\mathrm{i}\theta})=\mathrm{e}^{\mathrm{i}m\theta}$，管壁边界条件为绕 $S$ 的和乐等于 $\mathrm{e}^{\mathrm{i}m\theta}$，等价于插入磁荷 $m$ 的狄拉克奇异性。二维 Ising 模型中的无序场是更小的例子：绕它一周自旋反号；在 $G=SU(2)$ 的四维理论中，$T(1,S)$ 的管壁边界条件就是单极子荷 $n=1$。}
\physdef{无序算子的典型是 't Hooft 线：它规定场的奇异性，而不是场的取值；这需要 $S$ 的法丛定向，因此依赖 $\star$。在二维共形场论中 disorder operator 对应改变边界条件的插入（如 Ising 模型的 disorder 场），这是「无序」一词的来源。}
\physex{四维 $SU(2)$ 理论中 $T(1,S)$ 的无序算子是磁单极子线的插入，其奇异解在 $S$ 附近就是 BPS 单极子；OPE 的次领头项由单极子冒泡（Hecke 修正空间紧化的边界）给出。$U(1)$ 理论中两个无序算子的 OPE 由磁荷相加给出。}
\bridge{对应规则：无序算子 $\leftrightarrow$ Hecke 修正（改写丛），其标签是 ${}^LG$ 的最高权；修正空间 $\mathcal{Y}({}^Lw;x)$ 就是仿射 Grassmann 中的轨道。原文见 \S6.2 p.~81 与 \S9.1。}
\intuit{一句话：无序算子不测量场，而是「要求场长成某个样子」，这正是 Hecke 修正在几何上做的事——在一点上重新粘一次丛。}
\pitfall{按字面理解会以为 disorder 是「无序、随机」的意思；其实它是「不由场泛函定义」的意思。也不能忽略法丛定向：没有定向时 't Hooft 算子的符号有歧义，其 OPE 也就无从谈起。}
\end{word}

\begin{word}{特征膜}{eigenbrane}{eigenbrane}{\S6.4, p.~95}
\mathdef{设 $X$ 是一个线算子，它作用在膜范畴上给出自函子。膜 $\mathcal{B}$ 称为 $X$ 的特征膜（eigenbrane），若存在有限维复向量空间 $V$，使得 \[ X\,\mathcal{B}\simeq \mathcal{B}\otimes V . \] 更精确地，$V$ 可以随参数（例如 $C$ 上的点 $p$）变化而成为一个平坦向量丛的纤维；对一族可交换的线算子，要求同时对每个算子成立，此时称联合特征膜。数学上这与 Hecke 特征层的条件 $H_V(\mathcal{F})\simeq V_E\boxtimes\mathcal{F}$ 是同一个条件，只是把层换成了膜。}
\mathex{$G=SU(2)$、$C=\mathbb{P}^1$、$X=W_2$：支撑在 $\mathcal{M}_H$ 光滑点 $x$ 的零膜 $\mathcal{O}_x$ 满足 $W_2(\mathcal{O}_x)\simeq\mathcal{O}_x\otimes\mathcal{E}_q(2)|_x$，其中 $\mathcal{E}_q(2)$ 是泛丛在二维表示中的纤维；特征值是一个二维向量空间，随 $q\in C$ 拼成 $C$ 上的平坦丛。}
\physdef{物理上，特征膜是「被线算子作用后只差一个向量空间」的膜，也就是 BPS 态的支撑：它的 Chan-Paton 层被线算子作用后只整体乘上一个向量空间，故其上的态空间是线算子的特征空间。这里的「本征值」是向量空间或层，不是数。}
\physex{$\mathcal{M}_H$ 上的零膜是电特征膜（对 Wilson 算子而言），支撑在 Hitchin 纤维上、带平坦 Chan-Paton 丛的膜是磁特征膜（对 't Hooft 算子而言）；两者由 S-对偶互换。$SU(2)$、$C=\mathbb{P}^1$ 时 $\mathcal{M}_H$ 只有一个点，电特征膜就是这个点上的零膜。}
\bridge{对应规则：特征膜 $\leftrightarrow$ Hecke 特征层；电特征膜（零膜，$\Psi=\infty$）$\leftrightarrow$ $\mathrm{Bun}_G$ 上的 Hecke 特征层（$D$-模）；磁特征膜（$\Psi=0$）$\leftrightarrow$ 它在 S-对偶下的像。原文见 \S6.4 p.~95 与 \S8。}
\intuit{一句话：特征膜就是「线算子作用在它上面不改变它、只乘上一个向量空间」的膜，这正是特征层定义在膜范畴中的翻译。}
\pitfall{按字面理解会把特征膜当成一个普通向量；其实本征值是向量空间或层，膜本身是一个对象（支撑加 Chan-Paton 层），而且「特征」是对一整族算子（所有表示、所有点）而言的。}
\end{word}

\begin{word}{电特征膜}{electric eigenbrane}{electric-eigenbrane}{\S6.4, p.~97}
\mathdef{电特征膜是对所有 Wilson 算子 $W_R(S)$ 的联合特征膜：对每个表示 $R$ 与点 $p\in C$，存在向量空间 $V_R$ 使 \[ W_R(S)\,\mathcal{B}\simeq\mathcal{B}\otimes V_R . \] 在拓扑扭化理论中，这等价于说膜的 Chan-Paton 层在 Wilson 算子作用下按泛丛的表示纤维 $\mathcal{E}_p(R)$ 变换；若 $\mathcal{B}$ 是零膜，则 $V_R=\mathcal{E}_p(R)|_x$，其中 $x$ 是零膜的支撑点。}
\mathex{取 $\mathcal{M}_H(G)$ 的光滑点 $x$，零膜 $\mathcal{O}_x$ 的 Chan-Paton 层是天穹层 $U_x$；限制到点上一切丛都平凡，故特征值就是固定向量空间 $\mathcal{E}_q(R)|_x$。当 $q$ 沿 $C$ 变化时，这些向量空间拼成 $C$ 上的平坦丛；在复结构 $J$ 下它就是泛平坦 $G_{\mathbb{C}}$-丛，即几何 Langlands 一侧的 Langlands 参数。}
\physdef{电特征膜是 Wilson 算子的特征膜，物理上位于 $\Psi=\infty$ 一侧。它由不同离散电场 $\mathbf{e}_0$ 的扭化零膜之和构成；从本体上看它就是一个零膜（0-膜），支撑在 $\mathcal{M}_H$ 的一个点上，因此是「最简单」的膜。}
\physex{$G=SU(2)$、$C=\mathbb{P}^1$：$\mathcal{M}_H$ 只有一个点，零膜就是电特征膜，其特征值由泛 Higgs 丛在该点的纤维给出。对电特征膜作 S-对偶，得到 $\mathcal{M}_H({}^LG)$ 中支撑在 Hitchin 纤维上、带秩一平坦 Chan-Paton 丛的磁特征膜。}
\bridge{对应规则：电特征膜 $\leftrightarrow$ Hecke 特征层（$\mathrm{Bun}_G$ 上的 $D$-模）。零膜在镜像下对应 $\mathrm{Bun}_G$ 上的天穹层，后者满足 $H_V(\mathcal{O}_E)\simeq V_E\boxtimes\mathcal{O}_E$，即显然的 Hecke 特征性质。原文见 \S6.4 p.~97 与 \S8（p.~100 起）。}
\intuit{一句话：零膜只「看见」一点处的纤维，所以任何线算子作用上去都只能乘一个向量空间——这就是它自动成为特征膜的原因。}
\pitfall{按字面理解会以为电特征膜是需要额外构造的特殊膜；其实最简单的零膜就是电特征膜。也不要以为它的特征值是数：特征值是泛丛在表示下的纤维（一个向量空间），且随 $q\in C$ 变成平坦丛的纤维。}
\end{word}

\begin{word}{磁特征膜}{magnetic eigenbrane}{magnetic-eigenbrane}{\S6.4, p.~97}
\mathdef{磁特征膜是对所有 't Hooft 算子 $T({}^LR,S)$ 的联合特征膜。具体模型：取 Hitchin 纤维化 $\pi:\mathcal{M}_H(G)\to B$ 的一根纤维 $\mathcal{F}=\pi^{-1}(b)$，并在其上取秩一平坦线丛 $\mathcal{L}$，则以 $(\mathcal{F},\mathcal{L})$ 为支撑的膜满足 \[ T({}^LR,S)\,\mathcal{B}\simeq\mathcal{B}\otimes V_{{}^LR}, \] 其中 $S$ 与 $C$ 的交点给出 $V_{{}^LR}$ 的平坦丛。这样的膜称为 type $\mathcal{F}$ 膜。}
\mathex{$G=SU(2)$、$C$ 亏格 $g\ge2$、$b\in B$ 为 $\det\phi$ 的 $4g-4$ 个零点互异的正则值：纤维 $\mathcal{F}_b$ 是对应二重谱曲线 $S\to C$ 的 Prym 的 torsor（复环面，维数 $3g-3$），其上的平坦线丛由 $S$ 上的线丛给出。最小例子 $C=\mathbb{P}^1$ 时 $\mathcal{M}_H$ 本身是一个点，唯一的磁特征膜就是这个点上的膜。}
\physdef{磁特征膜是 't Hooft 算子的联合特征膜，物理上位于 $\Psi=0$ 一侧。其支撑是 Hitchin 纤维，这正是 \S10.5 直接从扩展 Bogomolny 方程读出的结论：$\varphi$ 的特征多项式跨过 't Hooft 算子不变，故 Hecke 变换把每根 Hitchin 纤维映到自身。}
\physex{$SU(2)$、$C$ 亏格 $2$：$\mathcal{F}_b$ 是维数 $3$ 的复环面；其上的平坦线丛给出 $T({}^LR,S)$ 的本征值。对磁特征膜作 S-对偶，得到电特征膜（零膜）；这就是两族特征膜互为 S-对偶像的含义。}
\bridge{对应规则：磁特征膜 $\leftrightarrow$ Hecke 特征层在 S-对偶下的像，其标签是 ${}^LG$ 的表示 ${}^LR$（即 Hecke 算子的标签）；支撑 $\mathcal{F}$ 与平坦丛一起对应 Langlands 参数。该纤维在辛形式 $\omega_K$ 下是特殊拉格朗日环面，故磁特征膜是拉格朗日膜；'t Hooft 算子的 OPE 系数正是这些膜交叠的态空间（相交上同调）的维数。原文见 \S6.4 p.~97 与 \S10.5 p.~182。}
\intuit{一句话：磁特征膜的支撑就是 Hitchin 纤维，而 't Hooft 算子在几何上只改写丛、不改 Higgs 场的特征多项式，所以整根纤维被保留下来。}
\pitfall{按字面理解会把磁特征膜与电特征膜当同一个膜的两个名字；它们分居 $\Psi=0$ 与 $\Psi=\infty$ 两侧，只由 S-对偶互换，并非同一个对象。也不要以为磁特征膜是某个单一向量：它是带平坦 Chan-Paton 丛的延展膜。}
\end{word}

\begin{word}{Hecke 特征层}{Hecke eigensheaf}{hecke-eigensheaf}{\S6.4, p.~95; \S10.4, p.~169}
\mathdef{设 $X$ 为紧黎曼面，$G$ 为约化群，${}^LG$ 为其 Langlands 对偶群，$\mathrm{Bun}_G$ 为 $X$ 上 $G$-丛的模叠。对 ${}^LG$ 的有限维表示 $V$ 与点 $x\in X$，Hecke 算子 $H_{V,x}$ 是 $D$-$模$范畴上的函子。$D$-模 $\mathcal{F}$ 称为 Hecke 特征层，若存在 $X$ 上的平坦 ${}^LG$-丛 $E$，使得对一切 $V$ 与 $x$ 有 \[ H_{V,x}(\mathcal{F})\simeq V_E\boxtimes\mathcal{F},\qquad V_E=E\times_{{}^LG}V . \] 当 $E$ 不可约时，这样的 $\mathcal{F}$ 唯一、不可约、全纯（holonomic 且具正则奇点）。}
\mathex{$G=GL_1$：$\mathrm{Bun}_{GL_1}=\mathrm{Pic}(X)$，$H_{1,x}$ 把 $\mathcal{L}$ 送到 $\mathcal{L}(x)$，特征层是 $\mathrm{Pic}$ 上 Fourier--Mukai 型的层。$G=GL_n$、$V=V_{\check\omega_1}$：Hecke 对应的纤维是 $\mathbb{P}^{n-1}$，$H$ 由沿 $\mathbb{P}^{n-1}$ 的推下给出；不可约 $GL_n$-局部系统 $E$ 给出特征层 $\mathcal{F}_E$。}
\physdef{物理上，Hecke 特征层是电特征膜的数学化身：在四维扭化规范理论中，Wilson 算子作用在膜上正如 Hecke 算子作用在 $D$-模上，而 Hecke 特征层正是「作用后只差一个 $V_E$」的对象。特征值是 $\mathrm{Bun}_G$ 上的向量丛 $V_E$，不是数。}
\physex{$G=SU(2)$、$C=\mathbb{P}^1$：$\mathrm{Bun}_{SU(2)}$ 只有一个稳定点，Hecke 算子把常数层映到 $H^\bullet(\mathbb{P}^1)\otimes\mathbb{C}$，即分次为 $0,2$ 的两维空间——常数层因此是（平凡的）Hecke 特征层。$G=SU(2)$、$C$ 亏格 $g\ge2$ 时特征层由 Hitchin 纤维上的平坦线丛给出。}
\bridge{对应规则：Hecke 特征层 $\leftrightarrow$ 特征膜。两者都满足「作用等于张量 $V_E$」这一条件，只是一个在 $\mathrm{Bun}_G$ 的 $D$-模范畴、一个在 $\mathcal{M}_H$ 的膜范畴；桥是镜对称（$B$-膜 $\leftrightarrow$ $A$-膜）加典范余各向同性膜。原文见 \S6.4 p.~95 与 \S10.4 p.~169。}
\intuit{一句话：Hecke 特征层就是「被所有 Hecke 算子同时对角化」的层，正如特征膜是被所有线算子同时对角化的膜；「对角化」在这里只能意味着乘一个固定的向量丛。}
\pitfall{按字面理解会把 Hecke 特征层当成「特征函数」或「Hecke 算子的特征向量」；实际上它是层或模，本征值是 $\mathrm{Bun}_G$ 上的向量丛 $V_E$，而且特征层的存在性正是几何 Langlands 的结论，不是定义的一部分。}
\end{word}

\begin{word}{Hecke 算子}{Hecke operator}{hecke-op}{\S9.1, p.~132}
\mathdef{设 $\mathrm{Hecke}_{V,x}\subset\mathrm{Bun}_G\times\mathrm{Bun}_G$ 为 Hecke 对应：$(E,E')$ 属于它，当且仅当 $E'$ 在 $x$ 外与 $E$ 同构、且在 $x$ 处的修正型落在 $V$ 的权轨道中（更精确地说，型由 $V$ 的支配最高权决定）。Hecke 算子是 \[ H_{V,x}(\mathcal{F})=p_{2*}\big(p_1^*\mathcal{F}\otimes\mathcal{L}_V\big), \] 其中 $p_1,p_2$ 是两个投影。它是 $D(\mathrm{Bun}_G)$ 上的自函子。}
\mathex{$G=GL_1$：对应为 $\{(\mathcal{L},\mathcal{L}(x))\}$，$H_{1,x}$ 是平移。$G=GL_2$、$V=V_{\check\omega_1}$：纤维为 $\mathbb{P}^1=\mathrm{Gr}(1,2)$；型 $(1,0)$ 的丛在 $x$ 外不变，在 $x$ 处把 $\mathcal{O}\oplus\mathcal{O}$ 改成 $\mathcal{O}(1)\oplus\mathcal{O}(-1)$。$G=SU(2)$、型 $n$：纤维为维数 $n$ 的仿射 Grassmann 轨道。}
\physdef{物理上 Hecke 算子是 't Hooft 线算子的作用：跨过 $x$ 处的 't Hooft 算子时丛被局部改写，'t Hooft 算子的标签 ${}^LR$ 正是 $V$ 的标签。在 $\Psi=\infty$ 一侧，Wilson 算子以张量积方式实现同一算子的 S-对偶版本。}
\physex{Hecke 算子的「乘法」是 OPE：$T({}^LR)T({}^LR')=T({}^LR+{}^LR')\oplus\cdots$，其代数与 ${}^LG$ 的张量积表示环同构（\S10.4）。$SU(2)$ 例：$T(w)T(w)=T(2w)\oplus T(0)$，对应 $\mathbf{2}\otimes\mathbf{2}=\mathbf{3}\oplus\mathbf{1}$；次领头项的系数来自单极子冒泡。}
\bridge{对应规则：Hecke 算子 $H_{V,x}$ $\leftrightarrow$ 't Hooft 线算子 $T(V,x)$（$\Psi=0$），也 $\leftrightarrow$ Wilson 线算子 $W(V,x)$（$\Psi=\infty$，S-对偶版本）。自由参数是表示 $V$ 与点 $x$；原文见 \S9.1 p.~132 与 \S10.4。}
\intuit{一句话：Hecke 算子就是「只改一点、其余不动」的丛改写，这恰好是 't Hooft 线跨过时对场做的操作。}
\pitfall{按字面理解会把 Hecke 算子当成普通积分算子（乘一个核）；其实它是函子，在对应上作拉回与推下，其「乘法」是 OPE，含来自模空间奇异性的非平凡系数。也要注意 Wilson 侧与 't Hooft 侧是同一 Hecke 算子的两个 S-对偶实现，不是一个算子的两个不同定义。}
\end{word}

\begin{word}{Hecke 修正}{Hecke modification}{hecke-modification}{\S9.1, p.~132; \S9.3, p.~145}
\mathdef{固定 $x\in C$、$G$-丛 $E$。$E$ 在 $x$ 处的 Hecke 修正是二元组 $(E',\beta)$，其中 $E'$ 是 $C$ 上 $G$-丛，$\beta:E'|_{C\setminus\{x\}}\xrightarrow{\sim}E|_{C\setminus\{x\}}$ 是同构；若 $\beta$ 可延拓过 $x$，则修正是平凡的。修正的型由 $\beta$ 在 $x$ 附近的极点阶决定，即由双陪集 \[ G(\mathcal{O})\backslash G(\mathcal{K})/G(\mathcal{O}),\qquad \mathcal{O}=\mathbb{C}[[z]],\ \mathcal{K}=\mathbb{C}((z)) \] 的元素决定；对单连通 $G$ 由支配余权 $\lambda\in X_*(T)^+$ 分类。}
\mathex{$G=GL_1$：型 $k\in\mathbb{Z}$ 唯一，$E'=E\otimes\mathcal{O}(x)^k$，局部截面为 $s(z)=z^{-k}f(z)$（$f$ 全纯），$k\ne0$ 时次数改变。$G=GL_2$、$E=\mathcal{O}\oplus\mathcal{O}$：型 $(1,0)$ 的修正空间是 $\mathbb{P}^1=\mathrm{Gr}(1,2)$，生成截面为 $(c/z)s_1(0)+(h_1,h_2)$；型 $(2,0)$ 的空间是 $\mathrm{WCP}^2(1,1,2)$ 挖一点，其紧化还要加入型 $(1,1)$ 的 $0$ 维点。}
\physdef{物理上 Hecke 修正是 't Hooft 算子跨过时丛发生的变化：$x$ 处「局部重粘」，$x$ 外完全同构。修正的型就是 't Hooft 算子的 GNO 磁荷；全部修正构成仿射 Grassmann 流形 $\mathrm{Gr}_G$ 中的有限维轨道 $\mathcal{Y}({}^Lw;x)$。}
\physex{$SU(2)$、$C=\mathbb{P}^1$、电荷 $n$：型 $n$ 的修正空间维数为 $n$；其紧化 $\overline{\mathcal{Y}(n)}$ 由 associated 型 $n-2,n-4,\dots$ 分层，每一层由一次单极子冒泡实现。修正空间的相交上同调给出 $T(n)$ 作用在膜上的态空间。}
\bridge{对应规则：修正的型 $\leftrightarrow$ 't Hooft 算子的 GNO 磁荷 $\leftrightarrow$ ${}^LG$ 的最高权；修正空间 $\leftrightarrow$ 仿射 Grassmann 的 Schubert 胞腔。原文见 \S9.1 p.~132 与 \S9.3 p.~145。}
\intuit{一句话：修正就是「在 $x$ 处换一种粘法」，$\beta$ 记录换法、型记录换的程度；这正是 't Hooft 线的奇异性所指定的操作。}
\pitfall{按字面理解会以为 Hecke 修正是「在 $x$ 处允许极点」；其实 $x$ 外必须严格同构，且 $\beta$ 是数据的一部分，不同的 $\beta$ 给出不同的修正。也要注意「型」不是丛的同构类（后者由 Grothendieck 分类给出），而是 $G(\mathcal{O})$-轨道，两者是不同的分层。}
\end{word}

\begin{word}{Hecke 对应}{Hecke correspondence}{hecke-correspondence}{\S9.3, p.~150}
\mathdef{Hecke 对应是 $\mathrm{Bun}_G\times\mathrm{Bun}_G$ 中的子叠 \[ \mathrm{Hecke}_{V,x}=\big\{(E,E'):\ E'\ \text{在}\ x\ \text{外与}\ E\ \text{同构，修正型落在}\ V\ \text{的权轨道中}\big\}, \] 连同两个投影 $p_1,p_2$ 到 $\mathrm{Bun}_G$。$\mathrm{Hecke}_{\check\lambda,x}$ 的纤维是型 $\check\lambda$ 的修正空间（有限维）。它是一个对应而非单值映射：同一 $E$ 可对应多个 $E'$，因此诱导的是函子而非函数。}
\mathex{$G=GL_2$、型 $\check\omega_1$：纤维为 $\mathbb{P}^1$。$G=SU(2)$、型 $n$：纤维为维数 $n$ 的仿射 Grassmann 轨道 $\mathcal{Y}(n)$。$G=GL_1$：对应是 $\mathrm{Pic}$ 上的平移图 $\{(\mathcal{L},\mathcal{L}(x))\}$，纤维是一点，故 $H_{1,x}$ 就是平移。}
\physdef{物理上 Hecke 对应是 't Hooft 线算子的几何化：它记录「跨过算子前」与「跨过算子后」的丛对，算子的作用就是对 $E_-$ 求和并把结果写到 $E_+$ 的位置。在膜的层面，对应是 $\mathcal{M}_H\times\mathcal{M}_H$ 中的一个复拉格朗日子流形 $Q$（\S10.5）。}
\physex{$G=SU(2)$、$C$ 亏格 $2$：$Q\subset\mathcal{M}_H\times\mathcal{M}_H$ 为中维复子流形；在 $\phi$ 正则半单的点上，$Q$ 的纤维是有限集，因为全纯 Higgs 场不能在一点被修改而保持全纯，修正空间的维数恰好被 $\phi$ 的全纯性条件消掉。}
\bridge{对应规则：Hecke 对应 $\leftrightarrow$ 't Hooft 算子的时空轨迹 $\mathbb{R}\times\{p\}$；$\mathcal{M}_H$ 侧的 $Q$ 就是该算子在四维理论中作用在 $A$-膜上的实现。原文见 \S9.3 p.~150 与 \S10.5 p.~182。}
\intuit{一句话：对应不是算子，而是「前后态的关系」；算子是在这个关系上作拉回-推下得到的函子。}
\pitfall{按字面理解会以为对应纤维中的点就是修正前的丛；其实 $p_1$ 给 $E$、$p_2$ 给 $E'$，纤维中的点（修正数据）才决定修正后的丛。也不要把 Hecke 对应与 Hecke 算子混同：前者是几何对象（关系），后者是由它诱导的函子。}
\end{word}

\par\medskip\noindent\textbf{第二组（F）：单极子、Bogomolny 方程与仿射 Grassmann 流形}\par

\begin{word}{单极子}{monopole}{monopole}{\S6.2, p.~82}
\mathdef{设 $W$ 为三维黎曼流形，$G$ 为紧李群，$A$ 为 $W$ 上 $G$-联络，$\varphi_0$ 为 $\mathrm{ad}$-值函数。BPS 单极子是满足 Bogomolny 方程 \[ F=\star\,\mathrm{d}\varphi_0 \] 的场对 $(A,\varphi_0)$，并要求 $|\varphi_0|\to v\ne0$（$v$ 为常数）与有限能量。其磁荷是绕无穷远球面的通量，取值在 $\pi_1(G)$ 的一个共轭类中。等价地，它是 't Hooft 算子的奇异解：$T(\rho,S)$ 附近 $A$ 的曲率具有 $F\sim(\mathrm{i}/2)\star_3\mathrm{d}(1/r)\rho'(0)$ 的奇异性。}
\mathex{$G=U(1)$：狄拉克单极子 $F=(\mathrm{i}m/2)\star_3\mathrm{d}(1/r)$，$\oint\mathrm{i}F/2\pi=m$，奇点处须挖去。$G=SU(2)$：'t Hooft--Polyakov 单极子，磁荷 $n\in\mathbb{Z}$，BPS 质量 $M=4\pi v|n|/g^2$；$n=1$ 的模空间维数为 $4$（位置 $3$ 加整体规范相位 $1$）。一般荷 $k$ 的单极子，其 Nahm 数据是取值于 $k\times k$ 矩阵的 Nahm 方程解，极点只出现在区间的两个端点（标准数据里个数恒为 $2$），留数取 $\mathfrak{su}(2)$ 的 $k$ 维不可约表示；荷由留数表示的维数而非极点个数给出，各极点位置给出单极子沿某一方向的坐标。这就是单极子模空间的仿射 Grassmann 描述（更细致的讨论见深水区「Nahm 极点与 ADHMN 构造」条）。}
\physdef{物理上单极子是带磁荷的孤立解或粒子态：在四维理论中它表现为 't Hooft 线算子附近的奇异解，而不是四维中可直接观测的粒子（原书 \S6.2 p.~82 特别强调这一点）。在三维规范理论中它是真正的粒子；BPS 单极子保持一半超对称。}
\physex{$SU(2)$ 理论中 't Hooft 线 $T(1,S)$ 附近就是单极子的场；两根 't Hooft 线靠近时，单极子可能缩成一点并被已有单极子吞掉，这就是单极子冒泡（\S10.2）。$U(1)$ 理论中单极子的磁荷可直接由绕 $S$ 的磁通量读出。}
\bridge{对应规则：单极子 $\leftrightarrow$ Hecke 算子的几何实现。单极子模空间即 Bogomolny 模空间 $Z$，它的仿射 Grassmann 描述（Nahm 极点给出极点数据）给出 Hecke 修正空间的胞腔分解。原文见 \S6.2 p.~82 与 \S9--\S10。}
\intuit{一句话：单极子是 't Hooft 线附近「磁荷长出来的样子」，它的模空间正是 Hecke 算子的态空间。}
\pitfall{按字面理解会把单极子当成 't Hooft 算子本身：单极子是算子的奇异解（或态），'t Hooft 线是支撑在曲线上的线算子。也不要把 BPS 单极子与瞬子混淆：前者是三维的（或在四维中作为 't Hooft 线附近的奇异解），后者是四维的。}
\end{word}

\begin{word}{单极子冒泡}{monopole bubbling}{monopole-bubbling}{\S10.1, p.~148; \S10.2, p.~155}
\mathdef{考虑 't Hooft 算子 $T({}^Lw)$ 的修正空间 $\mathcal{Y}({}^Lw;x)$。它是仿射 Grassmann 流形中的有限维轨道，一般非紧；其紧化 $\overline{\mathcal{Y}({}^Lw)}$ 新增的边界点由「一单位磁荷缩成一点并被已有单极子吞掉」给出，这一退化称为单极子冒泡。用权语言：当 $m_i\ge m_j+2$ 时，$(m_i,m_j)\to(m_i-1,m_j+1)$，即一个权退化、两个权合并。}
\mathex{$G=SU(2)$、型 $n$：$\mathcal{Y}(n)$ 的维数为 $n$，紧化后加入型 $n-2,n-4,\dots$ 的层，每次加入对应一次冒泡。$G=U(2)$、${}^Lw=(1,0)$：$\mathcal{Y}\cong\mathbb{P}^1$ 的边界含型 $(0,0)$ 的点（磁荷被吞）。$GL_2$、型 $(2,0)$：$\overline{\mathcal{Y}((2,0))}=\mathrm{WCP}^2(1,1,2)$ 加入型 $(1,1)$ 的点。}
\physdef{物理上，单极子冒泡是瞬子缩成一点并入已有单极子的过程：ADHM 图像中 $\mathbb{R}^4$ 上瞬子的尺度模趋于零；在 Hecke 一侧就是修正空间的边界。它不是量子隧穿，而是尺度模的经典退化。}
\physex{$SU(2)$ 理论中 $T(w)T(w)=T(2w)\oplus T(0)$ 的次领头项系数来自冒泡：$\mathbb{P}^1$ 丛上四次上同调中收敛到奇点的那个形式给出平凡项系数 $1$；局部模型是 $A_1$ 奇点 $\mathbb{C}^2/\mathbb{Z}_2$ 的 Eguchi--Hanson 解消 $T^*\mathbb{P}^1$，其上唯一的 $L^2$ 调和 $2$-形式给出这个系数。}
\bridge{对应规则：冒泡 $\leftrightarrow$ associated 权（Schubert 循环的边界层），即 Hecke 修正空间紧化的边界，也是 OPE 系数的几何来源。原文见 \S10.1 p.~148 与 \S10.2 p.~155。}
\intuit{一句话：冒泡就是「小单极子被大单极子吃掉」，等价于权从 $(m+1,m-1)$ 合并为 $(m,m)$，所以它给出 Schubert 胞腔闭包中的一层。}
\pitfall{按字面理解会把冒泡当成量子过程（隧穿）；其实是经典模空间的边界退化（尺度模趋零）。也不要把「冒泡」与「瞬子」混同：瞬子只是 ADHM 一侧的图像，冒泡发生在 't Hooft 算子的修正空间上。}
\end{word}

\begin{word}{Bogomolny 方程}{Bogomolny equations}{bogomolny-equations}{\S9, p.~124}
\mathdef{设 $C$ 为黎曼面，$W=I\times C$，$A$ 为 $W$ 上 $G$-联络，$\varphi_0$ 为 $\mathrm{ad}$-值函数。Bogomolny 方程是 \[ F=\star\,\mathrm{d}\varphi_0, \] 其中 $\star$ 为三维 Hodge 星、$F$ 为 $A$ 的曲率；它是 $\mathbb{R}\times C$ 上自对偶方程在时间无关与时间反演不变两个假设下的约化，也可写成 $F-\phi\wedge\phi+\star\mathrm{d}\varphi_0=0$。}
\mathex{$G=SU(2)$、$C=\mathbb{P}^1$、$W=I\times\mathbb{P}^1$：方程的解在 $I$ 的两端分别趋于固定的 $A_-$ 与 $\varphi_0=0$，其模空间是型 $n$ 的 Hecke 修正空间（维数 $n$）。$G=U(1)$：$F=\star\mathrm{d}\varphi_0$ 的解由 $C$ 上的全纯数据给出，模空间是复环面。}
\physdef{Bogomolny 方程是 $\Psi=0$ 扭化理论的 BPS 方程：把 't Hooft 线算子插在 $W$ 的内点 $p_i$ 处，方程的解空间就是超对称构型空间 $Z$，其相交上同调是 Hecke 算子的态空间。它是 Hitchin 方程在 $\mathbb{R}\times C$ 上的版本，多出的 $y$ 方向给出算子的时间演化。}
\physex{$SU(2)$、$C=\mathbb{P}^1$、型 $n$：解模空间维数为 $n$；$n=1$ 时它是 $\mathbb{P}^1$，正是 $SU(2)$ 仿射 Grassmann 的胞腔。$U(1)$ 情形还原为狄拉克单极子的经典解。方程的无算子情形维数为 $0$，这与「无算子时构型空间是单点」一致。}
\bridge{对应规则：Bogomolny 方程的解空间 $\leftrightarrow$ Hecke 对应 $Q$ 的纤维 $\leftrightarrow$ Schubert 胞腔；解空间的相交上同调 $\leftrightarrow$ OPE 的多重度空间。原文见 \S9 p.~124 与 \S10.1。}
\intuit{一句话：Bogomolny 方程就是「把 Hitchin 方程沿时间方向抬高一层」，所以它的解模空间自然给出 Hecke 修正。}
\pitfall{按字面理解会以为 Bogomolny 方程只是物理的 BPS 条件；其实它就是 $\mathbb{R}\times C$ 上的 Hitchin 方程（自对偶方程的时间无关约化）。也不要把它与三维的 BPS 单极子方程混淆：这里是四维理论约化到三维的方程，两者形式相同但所在空间不同。}
\end{word}

\begin{word}{仿射 Grassmann 流形}{affine Grassmannian}{affine-grassmannian}{\S9.1, p.~132; \S9.3, p.~145}
\mathdef{仿射 Grassmann 流形是 \[ \mathrm{Gr}_G=G_{\mathbb{C}}((z))\big/G_{\mathbb{C}}[[z]], \] 其中 $G_{\mathbb{C}}[[z]]$ 为形式幂级数、$G_{\mathbb{C}}((z))$ 为形式洛朗级数。它的点可解释为 $C$ 上 $G$-丛在 $x$ 处的一次修正（模去两侧平凡化）。$G=GL_N$ 时 $\mathrm{Gr}_N=GL_N(\mathbb{C}((z)))/GL_N(\mathbb{C}[[z]])$。$G(\mathcal{O})$ 的作用把 $\mathrm{Gr}_G$ 分成有限维轨道 $\mathcal{Y}({}^Lw)$，其闭包是 Schubert 循环。}
\mathex{$G=SU(2)$：轨道由 $n\ge0$ 标号，$\mathcal{Y}(n)$ 是 $z^n$ 所在轨道，维数 $n$；$n=1$ 时它等于 $\mathbb{P}^1$，$n=2$ 时维数为 $2$。$G=GL_2$：$\mathcal{Y}((1,0))\cong\mathbb{P}^1$，$\mathcal{Y}((2,0))$ 是 $\mathrm{WCP}^2(1,1,2)$ 挖去一点。}
\physdef{物理上，仿射 Grassmann 流形是 't Hooft 算子修正空间的载体：型 ${}^Lw$ 的修正空间就是 $\mathrm{Gr}_G$ 中的轨道 $\mathcal{Y}({}^Lw;x)=G(\mathcal{O})\cdot\mathrm{diag}(z^{-m_1},\dots,z^{-m_N})$；$\mathrm{Gr}_G$ 的一个点就是「在 $x$ 处把丛重粘一次」的完整数据。}
\physex{$G=SU(2)$、电荷 $n$：$\mathcal{Y}(n)\subset\mathrm{Gr}_{SU(2)}$ 维数为 $n$；修正空间紧化时边界出现型 $n-2$ 等，由单极子冒泡实现。$G=PSU(2)$ 时 $m_i$ 可取半整数，修正空间落在扭化的仿射 Grassmann 版本中。}
\bridge{对应规则：$\mathrm{Gr}_G$ 的点 $\leftrightarrow$ Hecke 修正的数据 $\leftrightarrow$ 't Hooft 奇异性的标签 $g(z)=\mathrm{diag}(z^{-m_i})=\rho_{\mathbb{C}}(z)$，后者又等于 GNO 磁荷的指数映射。原文见 \S9.1 p.~132 与 \S9.3 p.~145。}
\intuit{一句话：仿射 Grassmann 流形就是「在同一点上把丛换一种粘法」的参数空间，因此它同时是修正空间与 't Hooft 标签的家园。}
\pitfall{按字面理解会把它与有限维 Grassmann 流形（子空间空间）混淆：仿射 Grassmann 是无穷维的，但按 $G(\mathcal{O})$-轨道分层后每一层有限维。也不要把它与 $\mathrm{Bun}_G$ 混淆：前者是局部数据，后者是整体对象。}
\end{word}

\begin{word}{Schubert 胞腔}{Schubert cell}{schubert-cell}{\S9.3, p.~147}
\mathdef{设 $\mathrm{Gr}_G=G_{\mathbb{C}}((z))/G_{\mathbb{C}}[[z]]$ 为仿射 Grassmann 流形。左乘 $G_{\mathbb{C}}[[z]]$ 把它分成有限维轨道；型 ${}^Lw$ 的 Schubert 胞腔是 \[ \mathcal{Y}({}^Lw)=G_{\mathbb{C}}[[z]]\cdot \mathrm{diag}(z^{-m_1},\dots,z^{-m_N}), \] 它同构于某个仿射空间，因而光滑、但是拟仿射的（不紧）。对 $G=GL_N$ 由支配权 ${}^Lw=(m_1\ge\cdots\ge m_N)$ 标号，其维数为 \[ \Delta=\sum_{i<j}(m_i-m_j). \]}
\mathex{$G=GL_2$：型 $(1,0)$ 的胞腔为 $\mathbb{P}^1$（维数 $1$）；型 $(2,0)$ 的胞腔维数为 $2$，即 $\mathrm{WCP}^2(1,1,2)$ 挖去一点。$G=U(N)$、基本表示：胞腔同构于 $\mathbb{P}^{N-1}$。$G=SU(2)$、型 $n$：胞腔维数为 $n$，$n=1$ 时就是 $\mathbb{P}^1$。}
\physdef{物理上 Schubert 胞腔就是 't Hooft 算子 $T({}^LR)$ 的修正空间 $\mathcal{Y}({}^Lw;x)$：固定算子的 GNO 磁荷后，所有可能的修正（$x$ 外同构）构成这个有限维光滑流形；它的上同调（等变、相交）就是算子作用在态空间上的代数数据。}
\physex{$G=SU(2)$、电荷 $n$：修正空间维数 $n$；$n=1$ 时是 $\mathbb{P}^1$，上同调分次为 $0,2$，对应 $T(1)$ 作用的两维态空间。$GL_2$、型 $(2,0)$：修正空间是三维加权射影空间挖一点，故不是紧的。}
\bridge{对应规则：Schubert 胞腔 $\leftrightarrow$ 固定磁荷的修正空间 $\leftrightarrow$ 一个 Hecke 修正的连续族；胞腔的维数 $\Delta$ 与单极子数的预期维数（指标）一致。原文见 \S9.3 p.~147。}
\intuit{一句话：胞腔是「磁荷固定后还能连续变化的全部粘法」；它是开的、光滑的，因为还没算上磁荷继续冒泡的边界。}
\pitfall{按字面理解会以为胞腔与丛同构类一一对应；其实胞腔按 $G(\mathcal{O})$-轨道分层，而与 Grothendieck 分层（按丛的同构类 $E\cong\bigoplus_i\mathcal{O}(x)^{k_i}$）是两套不同的分层，后者无穷维。也不要把胞腔与它的闭包混为一谈。}
\end{word}

\begin{word}{Schubert 循环}{Schubert cycle}{schubert-cycle}{\S9.3, p.~148}
\mathdef{型 ${}^Lw$ 的 Schubert 循环是 Schubert 胞腔的闭包 \[ \overline{\mathcal{Y}({}^Lw)}\subset\mathrm{Gr}_G . \] 它一般奇异，并分解为更低维胞腔的不交并，边界上出现 associated 权型：${}^Lw'$ 称为 ${}^Lw$ 的 associated 权，若 ${}^Lw'$ 是 ${}^Lw$ 作为 ${}^LG$ 的权但不是最高权，等价地 ${}^Lw-{}^Lw'$ 是伴随表示的非零支配权。$U(N)$ 情形：当 $m_i\ge m_j+2$ 时 $(m_i,m_j)\to(m_i-1,m_j+1)$ 得到 associated 权。}
\mathex{$G=GL_2$、型 $(1,0)$：胞腔已是紧的 $\mathbb{P}^1$，循环等于胞腔本身。型 $(2,0)$：循环为 $\mathrm{WCP}^2(1,1,2)$，其中含型 $(1,1)$ 的点，在该点处奇异。$G=SU(2)$、型 $n$：循环中加入型 $n-2,n-4,\dots$ 的层，逐层由单极子冒泡实现。}
\physdef{物理上 Schubert 循环是 't Hooft 算子复合时的次领头项：$T({}^Lw)T({}^Lw')$ 的 OPE 中，除最高项 $T({}^Lw+{}^Lw')$ 外，其余项由 associated 型给出，其系数是 Hecke 修正空间奇异性的产物（\S10.4）。循环的边界就是冒泡发生的方向。}
\physex{$G=SU(2)$：$T(w)T(w)=T(2w)\oplus T(0)$；次项 $T(0)$ 对应 $\overline{\mathcal{Y}(2)}$ 中型 $0$ 的边界层，其系数 $1$ 来自 $A_1$ 奇点解消上唯一的 $L^2$ 调和 $2$-形式。$U(2)$、型 $(1,0)$：闭包是 $\mathbb{P}^1$，无低维边界，故 $T(1)^2=T(2)$ 无次项。}
\bridge{对应规则：Schubert 循环 $\leftrightarrow$ associated 权型 $\leftrightarrow$ 修正空间紧化的边界层 $\leftrightarrow$ OPE 的次领头项与冒泡。原文见 \S9.3 p.~148 与 \S10.2。}
\intuit{一句话：循环是「胞腔加上它退化到的一切」，退化正是冒泡；算子代数正是靠这些边界才不是简单的加法。}
\pitfall{按字面理解会以为 Schubert 循环是光滑流形；其实闭包一般奇异，必须用相交上同调（$\mathrm{IC}$ 层）而不是普通上同调。也不要把循环与胞腔当成同一个对象：胞腔光滑开、循环奇异闭。}
\end{word}

\begin{word}{泛丛}{universal bundle}{universal-bundle}{\S7.2, p.~106}
\mathdef{设 $\mathcal{M}=\mathcal{M}(G,C)$ 为 $C$ 上稳定 $G$-丛的模空间。泛丛是 $\mathcal{M}\times C$ 上的 $G$-丛 $\mathcal{E}$，使得对每个 $p\in\mathcal{M}$，限制 $\mathcal{E}|_{\{p\}\times C}$ 同构于 $p$ 所代表的丛 $E_p$。对伴随群 $G_{\mathrm{ad}}$ 泛丛总存在；对单连通覆盖 $\bar G$，它一般不存在，障碍是 \[ \zeta=\xi(\mathcal{E}_{\mathrm{ad}})\in H^2(\mathcal{M},\mathbb{Z}), \] 但作为扭化向量丛总存在。与之配套的还有泛 Higgs 丛 $(\mathcal{E},\hat\varphi)$，其中 $\hat\varphi\in H^0(\mathcal{M}_H(G)\times C,\mathrm{ad}(\mathcal{E})\otimes K_C)$。}
\mathex{$G=SU(2)$、$C=\mathbb{P}^1$：$\mathcal{M}_H$ 为一点，泛丛退化为 $C$ 上的平凡丛。$G=SU(2)$、$C$ 亏格 $g\ge2$：$\dim\mathcal{M}_H=6g-6$，泛丛存在（$SU(2)$ 的障碍可消失），其限制 $E_p$ 随 $p$ 变动。$G=PSU(2)=SO(3)$：泛 $SO(3)$-丛总存在，但提升为 $SU(2)$-丛有 $\mathbb{Z}_2$ 障碍。}
\physdef{物理上泛丛是 Wilson 算子的载体：Wilson 算子的作用是把膜与 $\mathcal{E}_p(R)$ 张量（$R$ 为 $G$ 的表示），同时把离散参数 $\mathbf{e}_0$ 平移（\S8.1）；因此泛丛的每根纤维就是 Wilson 线在被穿过的点上的「特征值丛」。在 $B$-模型一侧，它就是 Hecke 特征层公式中的 $V_E$。}
\physex{$SU(2)$、$C=\mathbb{P}^1$：$\mathcal{M}_H$ 为一点，零膜的 Chan--Paton 层限制到该点时为固定向量空间 $\mathcal{E}_q(R)|_x$；$q$ 沿 $C$ 变动时它成为 $C$ 上的平坦丛，即 Langlands 参数。$g\ge2$ 时 $\mathcal{E}_p(R)$ 一般非平凡，这使 Wilson 算子真正改变膜。}
\bridge{对应规则：泛丛 $\mathcal{E}$ $\leftrightarrow$ Hecke 特征层公式中的 $V_E=E\times_{{}^LG}V$ $\leftrightarrow$ 磁特征膜上的平坦 Chan--Paton 丛；障碍 $\zeta$ 的平移 $\leftrightarrow$ Wilson 算子对 $\mathbf{e}_0$ 的平移。原文见 \S7.2 p.~106 与 \S8.1。}
\intuit{一句话：泛丛把「模空间每一点对应的丛」拼成一个整体对象，于是「逐点取丛」这句话才能写成几何操作。}
\pitfall{按字面理解会以为泛丛是 $\mathcal{M}$ 上的丛；其实它是 $\mathcal{M}\times C$ 上的丛，$\mathcal{M}$ 上的丛是它的推下/扭化版本。也不要在单连通群上默认它存在：$SU(N)$、$N\ge2$、$g\ge2$ 时它作为普通丛一般不存在，必须用扭化丛或 projective 版本。}
\end{word}

\begin{word}{超多重态}{hypermultiplet}{hypermultiplet}{\S4.1, p.~46}
\mathdef{四维 $\mathcal{N}=2$ 超对称代数中，超多重态（hypermultiplet）是超荷的一个不可约表示：包含两个 Weyl 费米子与四个实标量（即一个四元数）。四个实标量在超凯勒流形的切空间上给出三个复结构 $I,J,K$，故其标量场取值于超凯勒流形。与之相对，$\mathcal{N}=2$ 矢量多重态含规范场 $A_\mu$、两个 Weyl 费米子与一个复标量。}
\mathex{$\mathcal{N}=4$ 理论可视为 $\mathcal{N}=2$ 矢量多重态加一个伴随超多重态：六个实标量中两个属于矢量多重态、四个属于超多重态。$G=SU(2)$、$C$ 为黎曼面：全部场构成无穷维超凯勒流形 $X$，规范群 $\mathcal{G}=\mathrm{Maps}(C,G)$ 作用，超凯勒商 $X/\!\!/\mathcal{G}$ 即 Hitchin 模空间 $\mathcal{M}_H$（\S4.1 p.~46 附近）。}
\physdef{物理上超多重态是标量模空间的建材：超多重态标量天然参数化超凯勒流形，矢量多重态的规范场通过矩映射 $X\to\mathfrak{g}^*$ 作用，取矩映射零点再商掉复规范群便得到低能 sigma 模型的靶空间。$\mathcal{N}=4$ 理论的超凯勒结构（$\mathcal{M}_H$ 的 $I,J,K$ 三个复结构）正来自这个构造。}
\physex{$G=SU(2)$、$C=\mathbb{P}^1$：$X$ 为联络与 Higgs 场的空间，$\mathcal{G}$ 的复化作用自由，商有限维（因为 Hitchin 方程椭圆）；商的维数为 $\dim\mathcal{M}_H=6g-6$，$g=0$ 时为 $0$。$A$-模型取辛形式 $\omega_K$、$B$-模型取复结构 $J$，两者由超凯勒结构 $I,J,K$ 决定。}
\bridge{对应规则：超多重态 $\leftrightarrow$ Higgs 场 $\varphi$（其四个实标量在扭化后组合成复 1-形式 $\varphi$）$\leftrightarrow$ Hitchin 模空间的超凯勒结构 $\leftrightarrow$ 特征膜所在的 $\mathcal{M}_H$。原文见 \S4.1 p.~46 与 \S6--\S7。}
\intuit{一句话：超多重态就是把「额外四个实标量」打包成四元数，于是标量空间自动是超凯勒流形，这正是镜像对称能作用在 $\mathcal{M}_H$ 上的原因。}
\pitfall{按字面理解会以为超多重态只是「又一种粒子」；其实它在扭化理论中的作用是提供 Higgs 场与超凯勒结构，从而提供 $A$-模型、$B$-模型两个复结构。也不要把超多重态与矢量多重态混同：前者无规范场、标量是四元数，后者含规范场、标量只有一个复数。}
\end{word}

\par\medskip\noindent\textbf{本轨道小结}\par
T5 的一条主线是：Wilson 算子沿点 $x$ 作用 $\leftrightarrow$ Hecke 算子 $H_{V,x}$；'t Hooft 算子沿 $x$ 作用 $\leftrightarrow$ 同一算子的 S-对偶实现；两者在 $\Psi=0$ 与 $\Psi=\infty$ 两个扭化点分别成为拓扑算子。其几何核心是 't Hooft 算子跨过时丛在 $x$ 处被 Hecke 修正，全部修正构成仿射 Grassmann 流形中的 Schubert 胞腔，胞腔的闭包（Schubert 循环）的边界由单极子冒泡给出，而其相交上同调给出 OPE 系数，于是 Hecke 算子代数同构于 ${}^LG$ 的表示环。膜一侧的对应物是特征膜：零膜（电特征膜）被 Wilson 算子映为张量 $\mathcal{E}_p(R)$ 后的膜，Hitchin 纤维上的 type $\mathcal{F}$ 膜（磁特征膜）被 't Hooft 算子保持，二者分别就是 Hecke 特征层在 $\mathrm{Bun}_G$ 上的 $D$-模版本的物理化身。
